\documentclass[11pt]{article}
\usepackage[margin=1in]{geometry}
\usepackage{amsmath,amssymb}
\usepackage{booktabs}
\usepackage{hyperref}
\usepackage{listings}
\usepackage{xcolor}
\lstset{basicstyle=\ttfamily\small,breaklines=true,frame=single}

\title{Independent Replication Report --- arXiv:1806.11463 \\
\large Bayesian Deep Learning on a Quantum Computer}
\author{OpenClaw agent (Ollie/Rick Stevens replication project)}
\date{2026-07-03 (LaTeX rendering 2026-07-06)}

\begin{document}
\maketitle

\section{Paper and scope}

Zhao, Pozas-Kerstjens, Rebentrost, Wittek (2019), arXiv:1806.11463v3.
The paper's thesis is that infinite-width feedforward neural nets are
equivalent to Gaussian processes (GPs), and that GP posterior inference
reduces to a single primitive --- the linear-system solve
$(K + \sigma_n^2 I)^{-1} v$. The authors therefore propose the HHL
(Harrow--Hassidim--Lloyd) quantum matrix-inversion algorithm as a drop-in
core for Bayesian deep-learning inference, with a claimed polynomial
asymptotic speedup.

The paper's \emph{empirical} contribution is a small-scale demonstration
of the HHL core on a specific $2\times 2$ matrix
$A = \tfrac{1}{2}\bigl(\begin{smallmatrix}3&1\\1&3\end{smallmatrix}\bigr)$
on Rigetti QVM (Fig.~1,2) and on IBMQX5 / Rigetti 8Q-Agave hardware
(Fig.~3), with the headline hardware fidelity
$F=0.78$ on IBMQX5.

\section{Claims replicated}
\begin{itemize}
  \item \textbf{C1} --- GP posterior mean/variance reduce to
    $A^{-1}v$. \emph{Replicated exactly.}
  \item \textbf{C2} --- Noiseless HHL on the paper's $2\times 2$ matrix
    produces $A^{-1}|b\rangle$. \emph{Replicated (fidelity 1.000000).}
  \item \textbf{C3} --- Smooth monotone fidelity decay with gate
    noise. \emph{Replicated qualitatively.}
  \item \textbf{C4} --- IBMQX5 hardware $F=0.78$. \emph{Out of scope
    (device retired); bracketed by our noisy sweep at 5--10\%
    per-gate depolarizing.}
  \item \textbf{C5} --- Resource counts. \emph{Verified structurally
    (our shallow variant uses 4 qubits by exploiting the
    Hadamard-diagonal structure of $A$; the paper's generic
    circuit uses 6).}
\end{itemize}

\section{Verdict}

\textbf{REPLICATED} for the paper's testable simulation-tier claims.
The quantum primitive delivers the same Bayesian posterior mean
($0.475$) and variance ($0.6725$) as the classical inversion on the
toy GP, and the noise sweep reproduces the qualitative shape of
Fig.~1(a). The hardware headline ($F=0.78$ on IBMQX5) cannot be
re-run --- IBMQX5 has been retired --- but our noisy simulator
brackets it at 5--10\% per-gate depolarizing error, which is
realistic for 2018-era superconducting hardware.

\section{Honest critique}

This replication is a \emph{simulation-tier} success. Several important
caveats belong on the record:

\begin{enumerate}
  \item \textbf{No new Bayesian NN training.} The paper's headline is a
    Bayesian-inference primitive; it does not train a full Bayesian
    deep net. We likewise do not retrain any neural network --- we
    verify the equivalence at the GP layer (the paper's own reduction).
    So the replication does \emph{not} independently verify any
    predictive-accuracy or calibration metric on a specific dataset;
    the paper does not report any such metric either.
  \item \textbf{Trivial-scale kernel matrix.} The end-to-end GP posterior
    is on a $2 \times 2$ kernel. This matches what the paper actually
    ran on hardware, so it is a fair replication of the paper's
    demonstration, but it says nothing about whether the HHL primitive
    is useful at scale. The known caveats on HHL --- state-preparation
    cost, condition-number scaling, tomography-cost blowup at readout,
    the Aaronson ``fine print'' --- are not exercised here and are
    not addressed by the paper's demonstration either.
  \item \textbf{Shallow-circuit shortcut.} We use a 4-qubit HHL variant
    that exploits the Hadamard-diagonal structure of the paper's
    specific $A$. This is fair for reproducing C2/C3 numerically (the
    paper's own hardware run used a similar problem-specific
    circuit, ref.~[49]), but it is not the generic HHL circuit and
    it does not stress-test QPE for larger matrices.
  \item \textbf{No baseline comparison.} No non-Bayesian baseline or
    non-HHL Bayesian baseline (e.g., modern SGLD, HMC, variational
    inference) is compared here. The paper's own comparison is
    likewise absent --- the paper argues for an asymptotic complexity
    win, not a wall-clock win against 2018 classical MCMC/VI.
  \item \textbf{Uncertainty calibration not tested empirically.} We
    reproduce the GP predictive variance value ($0.6725$), but we do
    not have a real test set to check whether this quoted uncertainty
    is well-calibrated (coverage of nominal credible intervals). The
    paper does not test this either --- again, the empirical work
    is a hardware-primitive demonstration, not a calibration study.
  \item \textbf{Hardware claim (C4) unreproducible.} IBMQX5 has been
    decommissioned. We bracket the reported number via simulation
    rather than re-running on hardware. A modern re-run on IBM
    Quantum's current superconducting hardware (e.g., Heron, Eagle)
    would be a strictly stronger check.
\end{enumerate}

\section{Evidence}
See \texttt{report/evidence/} for JSON artifacts (noiseless HHL
statevector, 8-point noisy sweep, GP posterior, circuit listing,
version manifest) and \texttt{code/} for the two Python drivers.

\end{document}
